% Part: first-order-logic
% Chapter: proof-systems
% Section: natural-deduction

\documentclass[../../../include/open-logic-section]{subfiles}

\begin{document}

\iftag{FOL}
      {\olfileid{fol}{prf}{ntd}}
      {\olfileid{pl}{prf}{ntd}}

\olsection{فطری استنباط}

فطری استنباط اک !!a{derivation} نظام اے جہدا مقصد اصل استدلال (خاص طور
تے ریاضی داناں دے ورتے ہوئے باقاعدہ استدلال) دی صورت نوں پیش کرنا اے۔
اصل استدلال کئی ``فطری'' نمونیاں راہیں اگے ودھدا اے۔ مثال لئی، صورتاں
راہیں ثبوت سانوں اک فصلی مقدمے دی بنیاد اُتے کوئی نتیجہ قائم کرن دی
اجازت دیندا اے، ایس گل نوں قائم کر کے کہ اوہ نتیجہ ہر اک فصلی جز توں
نکلدا اے۔ بالواسطہ ثبوت سانوں کوئی نتیجہ ایس گل نال قائم کرن دی اجازت
دیندا اے کہ اوہدی نفی اک تضاد تک لے جاندی اے۔ شرطیہ ثبوت ``جے \dots
تاں \dots'' شرطیہ دعوا ایس گل نال قائم کردا اے کہ تالی، مقدم توں نکلدا
اے۔ فطری استنباط ایہناں فطری استنباطاں وچوں کجھ دی رسمی تشکیل اے۔ ہر منطقی
رابطے تے مسوّر دے نال دو قاعدے، اک تعارفی تے اک اخراجی قاعدہ، آؤندے
نیں تے ہر قاعدہ ایہو جہے کسے فطری استنباطی نمونے دے مطابق ہُندا اے۔
مثال لئی، $\Intro{\lif}$ شرطیہ ثبوت دے مطابق اے تے $\Elim{\lor}$
صورتاں راہیں ثبوت دے۔ اک خاص طور تے سادہ قاعدہ $\Elim{\land}$ اے،
جیہڑا $!A \land !B$ توں~$!A$ (یا $!B$) دا استنباط کرن دی اجازت دیندا
اے۔

فطری استنباط نوں ہور !!{derivation} نظاماں توں وکھ کرن والی اک خاصیت
اوس ولوں مفروضیاں دی ورتوں اے۔ فطری استنباط وچ !!^a{derivation}،
!!{formula}s دا اک درخت ہُندا اے۔ اک !!{formula}، !!{formula}s دے درخت
دی جڑ اُتے کھلوتا ہُندا اے، تے درخت دے ``پتے'' اوہ !!{formula}s ہُندے
نیں جنہاں توں نتیجہ اخذ کیتا جاندا اے۔ فطری استنباط وچ کجھ پتے والے
!!{formula}s، !!{derivation} دے اندر کردار نبھاؤندے نیں پر جدوں
!!{derivation} نتیجے تک پہنچدا اے تاں اوہ ``ورت کے مکائے'' جا چکے ہُندے
نیں۔ ایہ اصل استدلال وچ ایہو جہے مفروضے متعارف کراؤن دے عمل دے مطابق
اے جیہڑے تھوڑھے چِر لئی ای لاگو رہندے نیں۔ مثال لئی، صورتاں راہیں ثبوت
وچ اسیں ہر فصلی جز نوں سچ منّ لیندے آں؛ شرطیہ ثبوت وچ اسیں مقدم نوں
سچ منّ لیندے آں؛ بالواسطہ ثبوت وچ اسیں نتیجے دی نفی نوں سچ منّ لیندے
آں۔ وقتی مفروضے متعارف کراؤن تے فیر درمیانی قدم قائم کرن لئی اوہناں
نوں ختم کر دین دا ایہ طریقہ فطری استنباط دی نمایاں خاصیت اے۔ فطری
استنباطی !!{derivation} دے پتیاں اُتے فارمولیاں نوں مفروضے آکھیا جاندا
اے، تے استنباط دے کجھ قاعدے اوہناں نوں ``!!{discharge}'' کر سکدے نیں۔
مثال لئی، جے ساڈے کول کجھ مفروضیاں توں~$!B$ دا !!a{derivation} اے تے
اوہناں مفروضیاں وچ~$!A$ وی شامل اے، تاں $\Intro{\lif}$ دا قاعدہ سانوں
~$!A \lif !B$ دا استنباط کرن تے~$!A$ دے روپ والے ہر مفروضے نوں
خارج کرن دی اجازت دیندا اے۔ (ایہ حساب رکھن لئی کہ کیہڑے
استنباط اُتے کیہڑے مفروضے خارج کیتے جاندے نیں، اسیں استنباط
تے اوس ولوں خارج کیتے مفروضیاں اُتے اک نمبر دا نشان لاؤندے
آں۔) جیہڑے مفروضے !!{derivation} دے آخر اُتے !!{undischarged} رہندے
نیں، اوہ رل کے نتیجے دے صدق لئی کافی ہُندے نیں، ایس لئی
!!a{derivation} قائم کردا اے کہ اوہدے !!{undischarged} مفروضیاں توں
اوہدا نتیجہ منطقی طور تے نکلدا اے۔

فطری استنباط اُتے مبنی تعلق $\Gamma \Proves !A$ اُدوں تے صرف اُدوں
قائم ہُندا اے جدوں کوئی !!a{derivation} ایہو جہا ہووے کہ $!A$ درخت دا
آخری !!{sentence} ہووے تے ہر !!{undischarged} پتا~$\Gamma$ وچ ہووے۔
$!A$ فطری استنباط وچ اُدوں تے صرف اُدوں اک قضیہ اے جدوں کوئی
!!a{derivation} ایہو جہا ہووے کہ $!A$ آخری !!{sentence} ہووے تے سارے
مفروضے !!{discharged} ہون۔ مثال لئی، ایہ !!a{derivation} وکھاؤندا اے
کہ $\Proves (!A
\land !B) \lif !A$:
\begin{prooftree}
  \AxiomC{$\Discharge{!A \land !B}{1}$}
  \RightLabel{\Elim{\land}}
  \UnaryInfC{$!A$}
  \DischargeRule{\Intro{\lif}}{1}
  \UnaryInfC{$(!A \land !B) \lif !A$}
\end{prooftree}
نشان~$1$ دسدا اے کہ مفروضہ $!A \land !B$ نوں \Intro{\lif}
استنباط اُتے !!{discharged} کیتا جاندا اے۔

فطری استنباط وچ سیٹ~$\Gamma$ اُدوں تے صرف اُدوں متضاد اے جدوں
$\Gamma \Proves \lfalse$۔ قاعدہ \FalseInt{} ایس گل نوں یقینی بناؤندا
اے کہ کسے متضاد سیٹ توں ہر !!{sentence} نوں !!{derive} کیتا جا سکدا اے۔

فطری استنباطی نظام Gerhard Gentzen تے Stanis\l{}aw Ja\'skowski ولوں
1930 دی دہائی وچ بنائے گئے، تے بعد وچ Dag Prawitz تے Frederic Fitch
ولوں ہور اگے ودھائے گئے۔ چونکہ ایہدے استنباط ثبوت دے فطری طریقیاں دی
صورت نوں پیش کردے نیں، ایس لئی فلسفی ایہنوں پسند کردے نیں۔ Fitch ولوں
تیار کیتے روپ اکثر منطق دیاں تعارفی درسی کتاباں وچ ورتے جاندے نیں۔
منطق دے فلسفے وچ، فطری استنباط دے قاعدیاں نوں کدی کدی منطقی رابطیاں دے
معنی دین والا منّیا گیا اے (``ثبوتی-نظریاتی معنیات'')۔

\end{document}
